\exercisesection{Complete intersections}

\begin{enumerate}
\def\labelenumi{\arabic{enumi})}
\item
  Let \(k\) a field, A the quotient ring of the polynomial ring \(k[X,Y,Z]\) by the ideal generated by \(Z^2,ZX,Z(Y-1)\) , \(x\) and \(y\) the classes of X and Y in \(\Lambda\) . Show that the sequence \((x,y)\) is completely secant for A but is not A-regular, although the sequence \((y,x)\) is A-regular.
\item
  Let A be a regular local ring, and I an ideal of A.
\end{enumerate}

\begin{enumerate}
\def\labelenumi{\alph{enumi})}
\tightlist
\item
  If ht(I) = 1, prove that the following conditions are equivalent:
\end{enumerate}

\begin{enumerate}
\def\labelenumi{(\roman{enumi})}
\item
  the ideal I is principal;
\item
  A/I is a Gorenstein ring;
\item
  A/I is a Macaulay ring.
\end{enumerate}

\begin{enumerate}
\def\labelenumi{\alph{enumi})}
\setcounter{enumi}{1}
\item
  Assume \(ht(I)=2\) for A/I to be a Gorenstein ring, it is necessary and sufficient that the ideal I be completely secant (if \((L,p)\) is a minimal projective resolution of I; observe that \(L_{1}\) is of rank 1 and consequently \(L_{0}\) of rank 2).
\item
  Let \(k\) a commutative field; we take \(A = k[[X,Y,Z]]\) and \(I = (XY,YZ,ZX, X^2 - Y^2, Y^2 - Z^2)\) . Prove that \(I\) is an ideal of height 3 that is not completely secant, and that \(A/I\) is a Gorenstein ring.
\end{enumerate}

\begin{enumerate}
\def\labelenumi{\arabic{enumi})}
\setcounter{enumi}{2}
\tightlist
\item
  Let R be a regular local ring, I an ideal of R, and A the ring R/I. Let \(\mathbf{x} = (x_{1}, \ldots, x_{d})\) a coordinate system of R, and \(\mathbf{y} = (y_{1}, \ldots, y_{m})\) a minimal generating system of I, so that \(m = [I/\mathfrak{m}_{R}\mathbf{l} : \kappa_{R}]\) .
\end{enumerate}

\begin{enumerate}
\def\labelenumi{\alph{enumi})}
\item
  When I is contained in \(\mathfrak{m}_{\mathrm{R}}^{2}\) define a canonical isomorphism of \(\mathrm{H}_1(\mathbf{x},\mathrm{A})\) on \(\mathrm{I} / \mathfrak{m}_{\mathrm{R}}\mathrm{I}\) (consider an exact sequence of Koszul complexes).
\item
  For the images of the \(x_{i}\) in A to form a minimal system of generators of \(m_{A}\) it is necessary and sufficient that 1 be contained in \(m_{R}^{2}\) .
\item
  Prove that the dimension of the \(\kappa_{\mathrm{A}}\) -vector space \(H_1(\mathbf{x}, A)\) is equal to \(m - d + [\mathfrak{m}_{\mathrm{A}} / \mathfrak{m}_{\mathrm{A}}^{2} : \kappa_{\mathrm{A}}]\) .
\end{enumerate}

4) Let A be a Noetherian local ring. Set \(\delta(A) = [\mathfrak{m}_{A}/\mathfrak{m}_{A}^{2}:\kappa_{A}] - \dim(A)\) (cf. § 4, no. 5).

\begin{enumerate}
\def\labelenumi{\alph{enumi})}
\item
  Show that the dimension of the \(\kappa_{\Lambda}\) -vector space \(H_{i}(\mathbf{z}, A)\) , where \(\mathbf{z}\) is a minimal generating system of \(\mathfrak{m}_{\mathrm{A}}\) is independent of the choice of \(\mathbf{z}\) ; it is denoted \(h_{i}(\mathrm{A})\) .
\item
  Prove the equality \(h_i(\mathrm{A}) = h_i(\widehat{\mathrm{A}})\) for every \(i\) .
\item
  Prove that one has \(h_1(A) \geqslant \delta(A)\) , and equality holds if and only if \(A\) is a complete intersection ring (Remark 2 of No. 2; reduce to the case where \(A\) is complete and apply Exerc. 4 by presenting \(A\) as a quotient \(R/I\) , where \(R\) is a regular local ring and \(I \subset m_R^2\) ).
\item
  Let \((x_{1},\ldots,x_{p})\) a completely secant sequence of elements of \(m_{A}\) generating an ideal J. Prove the equality \(h_{1}(A/J)-\delta(A/J)=h_{1}(A)-\delta(A)\) (reduce to the case p=1 and distinguish two cases, according as \(x_{1}\) belongs or does not belong to \(m_{A}^{2}\) ).
\end{enumerate}

\begin{enumerate}
\def\labelenumi{\arabic{enumi})}
\setcounter{enumi}{4}
\tightlist
\item
  Let R be a regular local ring, I an ideal of R, A the ring R/I. For A to be a complete intersection ring, it is necessary and sufficient that the ideal I be completely secant (use Exercises 4 c) and 3 c)).
\end{enumerate}

¶ 6) Let A be a Noetherian local ring, I an ideal of A.

\begin{enumerate}
\def\labelenumi{\alph{enumi})}
\tightlist
\item
  Prove that the following conditions are equivalent:
\end{enumerate}

\begin{enumerate}
\def\labelenumi{(\roman{enumi})}
\item
  the ideal I is completely secant;
\item
  the A/I-module I/I² is free and one has dp\_A(A/I) \textless{} +∞.
\end{enumerate}

Under hypothesis (ii), deduce from exerc. 2 of § 3 that there exists a cancellable element \(x\) in I such that \(x \notin \mathfrak{m}_{\mathrm{A}}\mathrm{I}\) ; reason by induction on the dimension of \(\mathrm{I} / \mathrm{I}^2\) .)

\begin{enumerate}
\def\labelenumi{\alph{enumi})}
\setcounter{enumi}{1}
\item
  Prove that the set of prime ideals p of A for which the ideal \(I_{p}\) is completely secant is open in Spec(A).
\item
  Let B be a presentable ring; prove that the set of prime ideals q of B such that \(B_{q}\) let a complete intersection ring be open in Spec(B).
\end{enumerate}

\begin{enumerate}
\def\labelenumi{\arabic{enumi})}
\setcounter{enumi}{6}
\tightlist
\item
  Let A be a ring, B an A-algebra; we assume that the A-module B is free of finite type. In the following exercises, we denote by B* the A-module Hom \(_{A}\) (B, A), and we endow it with its natural structure as a B-module.
\end{enumerate}

\begin{enumerate}
\def\labelenumi{\alph{enumi})}
\item
  Let \((e_i)_{i\in I}\) a basis of the A-module B, and \((e_i^*)_{i\in I}\) the dual basis of \(\mathbf{B}^*\) Prove that the element \(\mathrm{Tr}_{\mathrm{B / A}}\) of \(\mathbf{B}^*\) (A, III, p.~110) is equal to \(\sum_{j}e_{i}e_{i}^{*}\) .
\item
  We assume that the B-module \(B^{*}\) is free of rank 1; let \(\ell\) an element forming a basis of this module, and let \(\delta\) the element of B such that \(Tr_{B/A} = \delta\ell\) . The ideal \(\delta B\) of B does not depend on the choice of \(\ell\) ; it is called the different ideal of B over A. Prove that \(N_{B/A}(\delta)\) is a generator of the discriminant ideal of B over A (A, III, p. 115; compute the discriminant of a basis ( \(e_{i}\) ) of B on \(\Lambda\) using the matrix of multiplication by \(\delta\) in this basis).
\end{enumerate}

\begin{enumerate}
\def\labelenumi{\arabic{enumi})}
\setcounter{enumi}{7}
\tightlist
\item
  Let A be a ring, P a monic polynomial of degree \(n\) of A{[}X{]}, let B be the A-algebra A{[}X{]}/(P). We denote by \(x\) the class of X in B. Let \(\ell: \mathrm{B} \to \mathrm{A}\) be the A-linear form such that \(\ell(x^{n-1}) = 1\) , \(\ell(x^i) = 0\) for \(0 \leqslant i \leqslant n-2\) ; we denote \(\tilde{\ell}: \mathrm{B}[\mathrm{X}] \to \mathrm{A}[\mathrm{X}]\) the A-map{[}X{]}-linear map deduced from \(\ell\) by extension of scalars.
\end{enumerate}

\begin{enumerate}
\def\labelenumi{\alph{enumi})}
\item
  Let \(\sigma : B[X] \to A[X]\) the map \(A[X]\) -linear such that \(\sigma(x^i) = X^i\) for \(i = 0, \ldots, n - 1\) . Prove the formula \(P\tilde{\ell}(Q) = \sigma((X - x)Q)\) for every element \(Q\) of \(B[X]\) (check it for \(Q = 1, \ldots, x^{n-1}\) ).
\item
  Let \(\mathrm{P}^{*}(\mathrm{X}) = b_{n-1}\mathrm{X}^{n-1} + \ldots + b_{0}\) the polynomial B{[}X{]} such that one has \(\mathrm{P}(\mathrm{X}) = (\mathrm{X} - x)\mathrm{P}^{*}(\mathrm{X})\) in B{[}X{]}. Prove that \((b_{0}\ell, \ldots, b_{n-1}\ell)\) is the basis of \(B^{*}\) dual to the basis \((1, \ldots, x^{n-1})\) of B (apply a) to polynomials \(x^{i}\mathrm{P}^{*}\) ). Deduce that the B-module \(B^{*}\) is free of rank one, and that \(\ell\) forms a basis.
\item
  Prove the equality. \(\mathrm{Tr}_{\mathrm{B / A}} = \mathrm{P}'(x)\ell\) in \(\mathrm{B}^*\) (cf.~exerc. 7 a)).
\end{enumerate}

\begin{enumerate}
\def\labelenumi{\arabic{enumi})}
\setcounter{enumi}{8}
\tightlist
\item
  Let A be a ring, n an integer, S the algebra \(A[X_{1},\ldots,X_{n}]\) (resp. \(A[[X_{1},\ldots,X_{n}]]\) ), and a an ideal of S generated by a sequence \(\mathbf{P}=(\mathbf{P}_{1},\ldots,\mathbf{P}_{n})\) completely secant for S. We set \(B=S/a\) and we denote by \(\pi\) the canonical homomorphism from S onto S/a. Assume that the A-module B is free of finite type. We propose to prove that the B-module \(B^{*}\) is free of rank 1 and the different ideal of B over A is generated by \(\pi(\det(\frac{\partial P_{i}}{\partial X_{j}}))\) .
\end{enumerate}

\begin{enumerate}
\def\labelenumi{\alph{enumi})}
\item
  We set \(C = S \otimes_{A} B\) and \(\xi_{i} = X_{i} \otimes 1 - 1 \otimes \pi(X_{i})\) for \(1 \leqslant i \leqslant n\) . We denote by \(\varphi : C \to B\) the homomorphism of A-algebras such that \(\varphi(P \otimes b) = \pi(P)b\) . Prove that the kernel of \(\varphi\) is generated by the sequence \(\boldsymbol{\xi} = (\xi_{1}, \ldots, \xi_{n})\) , and that this sequence is completely secant for C.
\item
  Let \((c_{ij}) \in \mathbf{M}_{n}(\mathbf{C})\) be a matrix satisfying \(P_{i} \otimes 1 = \sum_{j} c_{ij} \xi_{j}\) . Prove the equality \(\pi(\frac{\partial P_{i}}{\partial X_{j}}) = \varphi(c_{ij})\) .
\item
  The matrix \((c_{ij})\) defines a morphism of complexes \(\mathbf{K}_{\bullet}(\mathbf{P},\mathrm{C})\to\mathbf{K}_{\bullet}(\boldsymbol{\xi},\mathrm{C})\) (A, X, p.~151), whence a morphism of complexes of S-modules \(u:\mathbf{K}_{\bullet}(\mathbf{P},\mathrm{S})\to\mathbf{K}_{\bullet}(\boldsymbol{\xi},\mathrm{C})\) The homomorphism \(u_{n}:S\to C\) maps \(1_{S}\) on the element \(d=\det(c_{ij})\) of C.
\item
  Let \(v: \mathrm{Hom}_{\mathrm{S}}(\mathrm{C}, \mathrm{S}) \otimes_{\mathrm{C}} \mathrm{B} \to \mathrm{B}\) be the B-linear homomorphism such that \(v(f \otimes 1) = \pi(f(d))\) ; prove that \(v\) is an isomorphism (identify \(v\) with \(\mathrm{H}^n(\mathrm{Homgr}_{\mathrm{S}}(u, 1_{\mathrm{S}}))\) , and observe that each of the complexes \(\mathbf{K}_{\bullet}(\mathbf{P}, \mathrm{S})\) and \(\mathbf{K}_{\bullet}(\boldsymbol{\xi}, \mathrm{C})\) defines a resolution of B by free S-modules of finite type).

\item
  We consider the composite map \(t: \mathrm{Hom}_{\mathrm{A}}(\mathrm{B}, \mathrm{A}) \to \mathrm{Hom}_{\mathrm{S}}(\mathrm{C}, \mathrm{S}) \to \mathrm{B}\) , where the first arrow is obtained by extension of scalars and the second is deduced from \(v\) . Prove that \(t\) is B-linear and bijective (identify \(\mathrm{Hom}_{\mathrm{S}}(\mathrm{C}, \mathrm{S})\) with \(\mathrm{Hom}_{\mathrm{A}}(\mathrm{B}, \mathrm{A}) \otimes_{\mathrm{B}} \mathrm{C}\) ).
\item
  Prove the formula \(t(\mathrm{Tr}_{\mathrm{B / A}}) = \varphi (d)\) (express \(d\) and \(\mathrm{Tr}_{\mathrm{B / A}}\) using a basis of B and the dual basis of \(\mathrm{B}^*\) cf. exercise 7, a)).
\item
  Deduce from b), e), and f) that the different ideal of B over A is generated by \(\pi(\det(\frac{\partial P_i}{\partial X_j}))\) .
\end{enumerate}

\begin{enumerate}
\def\labelenumi{\arabic{enumi})}
\setcounter{enumi}{9}
\tightlist
\item
  Let A be a discrete valuation ring, v its normalized valuation, and K its field of fractions. Consider a Noetherian A-algebra B equipped with a homomorphism of A-algebras \(\varepsilon:B\to A\) Let I denote the kernel of \(\varepsilon\) and J its annihilator in B.
\end{enumerate}

\begin{enumerate}
\def\labelenumi{\alph{enumi})}
\tightlist
\item
  Prove that the following conditions are equivalent:
\end{enumerate}

\begin{enumerate}
\def\labelenumi{(\roman{enumi})}
\item
  the \(\Lambda\) -module \(1 / I^2\) is of finite length;
\item
  the local ring of \(K \otimes_{A} B\) at the maximal ideal \(K \otimes_{A} I\) is isomorphic (as a K-algebra) to K ;
\item
  ; we have \(\mathrm{K}\otimes_{\mathrm{A}}\mathrm{B} = (\mathrm{K}\otimes_{\mathrm{A}}\mathrm{I})\oplus (\mathrm{K}\otimes_{\mathrm{A}}\mathrm{J})\)
\end{enumerate}

We now assume that these conditions are satisfied; we denote by \(\gamma(B, \varepsilon)\) , or simply \(\gamma(B)\) the length of the A-module \(I/I^{2}\) , and by \(\eta(B, \varepsilon)\) , or simply \(\eta(B)\) , that of the A-module \(A/\varepsilon(J)\) .

\begin{enumerate}
\def\labelenumi{\alph{enumi})}
\setcounter{enumi}{1}
\item
  Let b be an ideal of B contained in I; one considers the A-algebra \(B' = B/b\) and the homomorphism \(\varepsilon': B' \to A\) deduced from \(\varepsilon\) . Prove that \(B'\) also satisfies the conditions of a), and that one has \(\gamma(B', \varepsilon') \leqslant \gamma(B, \varepsilon)\) and \(\eta(B', \varepsilon') \leqslant \eta(B, \varepsilon)\) .
\item
  Assume that the A-module B is free of finite type. Prove that J is a free A-module of rank one, a direct factor in B, and that one has \(\mathrm{Tr}_{\mathrm{B / A}}(x) = \varepsilon (x)\) for every \(x\in J\)
\item
  Assume moreover that B is a Gorenstein ring; denote by \(\ell\) a basis of the B-module B* (§ 3, exerc. 5), and by δ a generator of the different ideal of B over A (exerc. 7, b)). Prove the equality η(B) = v(ε(δ)) (prove that \(\ell\)(J) is equal to A, and use c)).
\item
  One places oneself in the situation of \(b\) ), with \(\mathfrak{b} \neq 0\) ; assume that the A-algebras B and B' are finite and flat, and that B is a Gorenstein ring. Prove that one has \(\eta(B') < \eta(B)\) (observe that the image of J in B/ \(\mathfrak{m}_{A}\) B is the socle of B/ \(\mathfrak{m}_{A}\) B, hence is contained in \(\mathfrak{b}/\mathfrak{m}_{A}\mathfrak{b}\) ).
\end{enumerate}

¶ 11) One keeps the notation of the preceding exercise; assume that the rings A and B are local and complete.

\begin{enumerate}
\def\labelenumi{\alph{enumi})}
\item
  Show that B is isomorphic to the quotient of a formal power series algebra \(S = A[[X_{1}, \ldots, X_{n}]]\) by an ideal a contained in \((X_{1}, \ldots, X_{n})\) ; one may take \(n = \dim_{\kappa_{A}}(\mathfrak{n}/\mathfrak{n}^{2})\) , where n is the image of \(m_{B}\) in \(B/m_{A}B\) .
\item
  Prove that the following conditions are equivalent:
\end{enumerate}

\begin{enumerate}
\def\labelenumi{(\roman{enumi})}
\item
  B is a finite flat A-algebra and a complete intersection ring;
\item
  the ideal a is generated by a sequence \((\mathrm{P}_{1},\ldots,\mathrm{P}_{n})\) completely secant for \(\kappa_{A}\otimes_{A}S\) .
\end{enumerate}

(cf. exerc. 5.)

\begin{enumerate}
\def\labelenumi{\alph{enumi})}
\setcounter{enumi}{2}
\item
  If these conditions are satisfied, prove that one has \(\eta(B) = \gamma(B) = v(\varepsilon(\det(\frac{\partial P_i}{\partial X_j})))\) (use exercises 9 and 10, \(d\) )).
\item
  Show that one can find a sequence \((f_{1},\ldots,f_{n})\) of elements of a completely secant for \(\kappa_{A}\otimes_{A}S\) such that the A-algebra \(\mathrm{B}^{\prime}=\mathrm{S}/(f_{1},\ldots,f_{n})\) satisfies \(\gamma(\mathrm{B}^{\prime})=\gamma(\mathrm{B})\) .
\end{enumerate}

(Deduce from the exact sequence \(\mathfrak{a}/\mathfrak{a}^{2}\to\Omega_{\mathrm{A}}(\mathrm{S})\otimes_{\mathrm{S}}\mathrm{B}\to\Omega_{\mathrm{A}}(\mathrm{B})\to0\) (A, III, p.~137) an exact sequence of A-modules \(\mathfrak{a}/\mathfrak{a}^{2}\otimes_{\mathrm{S}}\mathrm{A}\to\Omega_{\mathrm{A}}(\mathrm{S})\otimes_{\mathrm{S}}\mathrm{A}\to\mathrm{I}/\mathrm{I}^{2}\to0\) ; take elements \(f_{1},\ldots,f_{n}\) of \(\mathfrak{a}\) whose images in \(\Omega_{\mathrm{A}}(\mathrm{S})\otimes_{\mathrm{S}}\Lambda\) form a basis over A of the image of \(\mathfrak{a}/\mathfrak{a}^{2}\otimes_{\mathrm{S}}\mathrm{A}\) .)

\begin{enumerate}
\def\labelenumi{\alph{enumi})}
\setcounter{enumi}{4}
\tightlist
\item
  Hence, using Exercise 10, e), we have \(\gamma(B) \geqslant \eta(B)\) and equality holds if and only if B is a complete intersection ring.
\end{enumerate}

